\chapter{Tienden en honderdsten}\label{ch:g4:decimals}

Tussen $3$ en $4$ zwijgen de hele getallen --- en toch kan een hardloper
$3$ meter en een beetje voorliggen. De eenheid in tien en daarna in
honderd stukken knippen geeft de \emph{decimale \omterm{def:g4:fractions:def}{breuken}}, met een
prachtige korte schrijfwijze: het \omterm{def:g4:decimals:point}{decimaalteken}. (Decimale getallen krijgen
hun eigen hoofdstuk in \cref{ch:g5:decimals}.)

\section{Decimale breuken}

\begin{definition}[Tienden en honderdsten]\label{def:g4:decimals:def}
Een eenheid in $10$ gelijke delen knippen geeft \emph{tienden},
$\frac{1}{10}$; elk tiende weer in $10$ delen knippen geeft
\emph{honderdsten}, $\frac{1}{100}$. Dus
\[
1 = \frac{10}{10} = \frac{100}{100},
\qquad
\frac{1}{10} = \frac{10}{100}.
\]
\end{definition}

\begin{omfigure}
\begin{tikzpicture}[scale=0.42]
  \draw[gray!50, thin] (0,0) grid (10,10);
  \foreach \i in {0,...,9} \fill[omDef!35] (3,\i+0.06)
    rectangle (3.94,\i+0.94);
  \fill[omThm!70] (6.05,4.05) rectangle (6.95,4.95);
  \draw[thick] (0,0) rectangle (10,10);
  \node[below, font=\small] at (5,-0.4)
    {het hele vierkant $= 1$;\ \ een kolom (blauw) $= \frac{1}{10}$;\ \
    een vakje (rood) $= \frac{1}{100}$};
\end{tikzpicture}

{\small Het honderdvierkant: tien kolommen van tien vakjes. Eén kolom is
een tiende, één vakje een honderdste.}
\end{omfigure}

\begin{example}\label{ex:g4:decimals:count}
Als je $3$ kolommen en $7$ losse vakjes van het honderdvierkant inkleurt,
is er
\[
\frac{3}{10} + \frac{7}{100} = \frac{30}{100} + \frac{7}{100}
= \frac{37}{100}
\]
van het vierkant gekleurd --- zevenendertig honderdsten.
\end{example}

\section{Het decimaalteken}

\begin{definition}[Decimale schrijfwijze]\label{def:g4:decimals:point}
Het getal $2 + \frac{3}{10} + \frac{7}{100}$ schrijf je als
\[
2.37
\]
--- het \emph{decimaalteken}\index{decimaalteken} scheidt het hele deel
($2$) van het decimale deel; het eerste cijfer erachter telt de tienden,
het tweede de honderdsten. Dit boek zet als decimaalteken een punt, zoals
internationaal gebruikelijk is; je spreekt het uit als ``komma'', en op
school schrijf je meestal $2{,}37$. Lezen: ``twee komma drie zeven'', of
``twee en zevenendertig honderdsten''.
\end{definition}

\begin{example}\label{ex:g4:decimals:convert}
\[
\frac{7}{10} = 0.7, \qquad
\frac{43}{100} = 0.43, \qquad
\frac{5}{100} = 0.05, \qquad
3 + \frac{4}{10} = 3.4 .
\]
Pas op bij $\frac{5}{100}$: vijf \emph{honderdsten} hebben een $0$ op de
plaats van de tienden nodig --- $0.05$, niet $0.5$.
\end{example}

\begin{omfigure}
\begin{tikzpicture}[scale=1.35]
  \draw[-Stealth] (-0.2,0) -- (8.4,0);
  \foreach \i in {0,10,20,30,40,50,60,70,80}
    \draw ({\i*0.1},-0.09) -- ({\i*0.1},0.09);
  \foreach \i in {0,...,80} \draw ({\i*0.1},-0.045) -- ({\i*0.1},0.045);
  \foreach \x/\l in {0/2, 1/{2.1}, 2/{2.2}, 3/{2.3}, 4/{2.4},
                     5/{2.5}, 6/{2.6}, 7/{2.7}, 8/{2.8}}
    \node[below, font=\footnotesize] at (\x,-0.1) {$\l$};
  \fill[omThm] (3.7,0) circle (1.7pt)
    node[above, font=\small] {$2.37$};
\end{tikzpicture}

{\small Inzoomen tussen $2$ en $2.8$: grote streepjes bij elk tiende,
kleine streepjes bij elk honderdste. Het getal $2.37$ ligt tussen $2.3$ en
$2.4$, bij het $7$de kleine streepje.}
\end{omfigure}

\begin{example}[Geld bestaat uit honderdsten]\label{ex:g4:decimals:money}
Eén cent is één honderdste euro: een prijs van $4.85$ betekent $4$ hele
euro's en $85$ honderdsten. Met geld gebruik je het \omterm{def:g4:decimals:point}{decimaalteken} al elke
dag: $10$ cent $= 0.10$; twee euro vijftig $= 2.50$; en $0.05$ is een
muntje van vijf cent, niet van vijftig!
\end{example}

\section{Vergelijken met tienden}

\begin{method}[Decimale getallen vergelijken (eerste kennismaking)]\label{met:g4:decimals:compare}
\begin{enumerate}
  \item Vergelijk eerst de hele delen: $3.2 > 2.9$.
  \item Zijn de hele delen gelijk, vergelijk dan de tienden en daarna de
        honderdsten: $2.37 < 2.5$, want $3$ tienden $< 5$ tienden.
  \item Beide met twee decimalen schrijven helpt: $2.5 = 2.50$, en
        $37 < 50$ honderdsten.
\end{enumerate}
De valstrik: $2.37$ heeft meer cijfers dan $2.5$, maar is \emph{kleiner}
--- meer cijfers betekent niet groter!
\end{method}

\begin{example}\label{ex:g4:decimals:compare}
Zet $0.4$, $0.35$ en $0.09$ op volgorde: in honderdsten zijn dat $0.40$,
$0.35$ en $0.09$; dus
\[
0.09 < 0.35 < 0.4 .
\]
\end{example}

\section{Oefeningen}

\begin{exercise}[$\star$]\label{exo:g4:decimals:1}
Hoeveel is er op een honderdvierkant gekleurd als je inkleurt: $5$
kolommen; $2$ kolommen en $3$ vakjes; $45$ vakjes? Antwoord met \omterm{def:g4:fractions:def}{breuken}.
\end{exercise}

\begin{exercise}[$\star$]\label{exo:g4:decimals:2}
Schrijf met een \omterm{def:g4:decimals:point}{decimaalteken}:
$\frac{9}{10}$; \quad $\frac{27}{100}$; \quad $\frac{3}{100}$; \quad
$5 + \frac{6}{10}$; \quad $\frac{60}{100}$.
\end{exercise}

\begin{exercise}[$\star$]\label{exo:g4:decimals:3}
Schrijf als \omterm{def:g4:fractions:def}{breuk} (tienden of honderdsten):
$0.3$; \quad $0.51$; \quad $0.07$; \quad $1.9$.
\end{exercise}

\begin{exercise}[$\star$]\label{exo:g4:decimals:4}
In $6.58$: wat telt het cijfer $5$? En het cijfer $8$? Splits het getal
zoals in \cref{def:g4:decimals:point}.
\end{exercise}

\begin{exercise}[$\star$]\label{exo:g4:decimals:5}
Zet op een getallenlijn met streepjes per tiende van $4$ tot $5$:
$4.2$; \ $4.75$ (tussen welke twee streepjes?); \ $4.9$.
\end{exercise}

\begin{exercise}[$\star$]\label{exo:g4:decimals:6}
Neem over en vul in met $<$, $>$ of $=$:
\[
0.6 \;?\; 0.60, \qquad
2.8 \;?\; 2.75, \qquad
0.09 \;?\; 0.1, \qquad
5.3 \;?\; 5.13 .
\]
\end{exercise}

\begin{exercise}[$\star$]\label{exo:g4:decimals:7}
Zet op volgorde, van klein naar groot: $1.05$; \ $1.5$; \ $0.95$; \
$1.45$.
\end{exercise}

\begin{exercise}[$\star$]\label{exo:g4:decimals:8}
Schrijf in decimale schrijfwijze: drie euro en vijf cent; twaalf euro en
vijftig cent; negenennegentig cent. Zet daarna de drie prijzen op
volgorde.
\end{exercise}

\begin{exercise}[$\star$]\label{exo:g4:decimals:9}
Welke decimale getallen met één cijfer achter het \omterm{def:g4:decimals:point}{decimaalteken} liggen
strikt tussen $3$ en $4$? Noem ze allemaal. Hoeveel zijn er?
\end{exercise}

\begin{exercise}[$\star\star$]\label{exo:g4:decimals:10}
Tom zegt: ``$0.25$ is groter dan $0.5$, want $25$ is groter dan $5$''. Laat
hem met het honderdvierkant zien dat hij zich vergist (hoeveel vakjes
kleurt elk getal in?).
\end{exercise}

\begin{exercise}[$\star\star$]\label{exo:g4:decimals:11}
De drie sprongen van een atleet waren $3.6$ m, $3.58$ m en $3.65$ m. Zet de
sprongen op volgorde, van kort naar lang. Welke twee schelen precies
$\frac{5}{100}$ meter?
\end{exercise}
